\documentclass[11pt]{article}

% --- Codificación e idioma ---
\usepackage[T1]{fontenc}
\usepackage[utf8]{inputenc}
\usepackage[spanish,english]{babel}

% --- Matemáticas ---
\usepackage{amsmath,amsthm,amssymb}

% --- Tablas y formato ---
\usepackage{float}
\usepackage{booktabs}
\usepackage[normalem]{ulem}
\usepackage[a4paper,margin=2.5cm]{geometry}
\usepackage[hidelinks]{hyperref}
\newcommand{\repo}{\href{https://github.com/Segiraldo0610/ProyectoFinal_Algoritmos}{ProyectoFinal\_Algoritmos}}

% --- Numeración ---
\numberwithin{equation}{section}
\numberwithin{figure}{section}
\numberwithin{table}{section}

% --- Entornos ---
\theoremstyle{definition}
\newtheorem*{definition}{Definición}
\newtheorem*{restriccion}{Restricción}

% --- Nombres en español ---
\addto\shorthandsspanish{\spanishdeactivate{~<>}}

\title{Algoritmo de asignación de horarios estudiantiles\\
\large Entrega final: Implementación y sustentación}
\author{Cesar Santiago Canro Cabra \and Samuel Enrique Giraldo Sabogal}
\date{\today}

\begin{document}
\selectlanguage{spanish}

\maketitle

\begin{abstract}
En este documento se presenta el análisis y el diseño formal del problema de
asignación de horarios estudiantiles en la Pontificia Universidad Javeriana,
así como la estrategia de datos propuesta para probar y evaluar el algoritmo
de solución. Se formalizan las entradas, las precondiciones, las restricciones
de validez, las preferencias de los estudiantes y la salida esperada.
A partir de esta definición se describe la implementación del algoritmo
de asignación, se evalúan sus resultados en escenarios reproducibles y se
presenta la interfaz utilizada para la sustentación.
\end{abstract}

\tableofcontents
\newpage

\part{Definición formal del problema}

\section{Análisis del problema}

Se dispone de un conjunto de estudiantes que deben construir su horario
académico a partir de las asignaturas que desean cursar. Cada asignatura puede
ofrecer uno o varios grupos, y cada grupo se encuentra caracterizado por un
profesor, una o varias franjas horarias, un salón y una capacidad máxima de
estudiantes.

El problema consiste en determinar, para cada estudiante, qué grupos deben ser
asignados de manera que el horario resultante sea válido. Una asignación se
considera válida cuando cumple simultáneamente las siguientes condiciones:

\begin{enumerate}
  \item un estudiante no puede estar asignado a dos grupos cuyas franjas
        horarias se intersecten;
  \item un estudiante puede pertenecer, como máximo, a un grupo por cada
        asignatura solicitada;
  \item la cantidad de estudiantes asignados a un grupo no puede superar la
        capacidad máxima de dicho grupo;
  \item un estudiante solamente puede ser asignado a grupos correspondientes
        a asignaturas que haya solicitado.
\end{enumerate}

Satisfacer estas restricciones determina la \emph{validez} de una asignación,
pero no necesariamente una única solución. Para un mismo estudiante pueden
existir varios horarios válidos. Por esta razón, además de garantizar la
validez, se desea contar con información que permita evaluar la calidad del
horario obtenido.

Los principales criterios de calidad considerados son la prioridad de las
asignaturas, la preferencia por determinados profesores, la organización del
tiempo libre, la cercanía entre salones de clases consecutivas y la posibilidad
de coincidir con determinados compañeros.

Por tanto, en el problema se distinguen dos tipos de condiciones:

\begin{itemize}
  \item \textbf{Restricciones de validez:} condiciones obligatorias que todo
        horario generado debe satisfacer.
  \item \textbf{Criterios de calidad:} características que permiten evaluar y
        comparar horarios que ya son válidos.
\end{itemize}

El problema puede resumirse como la búsqueda de una asignación de grupos para
cada estudiante que produzca horarios válidos y cuya calidad pueda ser
evaluada mediante las prioridades y preferencias recibidas como entrada.
El análisis determina qué representa una solución correcta y qué información
resulta relevante para evaluarla. El procedimiento para obtenerla se describe
en la parte de implementación.

\section{Diseño del problema}

El diseño establece el contrato formal del problema con independencia del
algoritmo que posteriormente se utilice. Dicho contrato se organiza en cuatro
componentes:

\begin{enumerate}
  \item \textbf{Entradas:} estudiantes, asignaturas, grupos, profesores,
        salones, franjas horarias, solicitudes, prioridades y preferencias.
  \item \textbf{Precondiciones:} propiedades que deben satisfacer los datos
        recibidos antes de realizar la asignación.
  \item \textbf{Salida:} función que asocia a cada estudiante el conjunto de
        grupos que conforma su horario.
  \item \textbf{Postcondiciones:} restricciones que debe cumplir la salida
        para ser considerada válida.
\end{enumerate}

Una vez verificada la validez de la salida, se calculan métricas de calidad que
permiten comparar distintas asignaciones válidas. Las secciones siguientes
desarrollan cada componente del contrato.

\section{Especificación formal de las entradas}

\subsection{Conjuntos}

Se definen los siguientes conjuntos finitos:

\begin{itemize}
  \item $E = \{e_1,e_2,\ldots,e_n\}$: conjunto de estudiantes.
  \item $A = \{a_1,a_2,\ldots,a_m\}$: conjunto de asignaturas disponibles.
  \item $G = \{g_1,g_2,\ldots,g_k\}$: conjunto de grupos o secciones.
  \item $P = \{p_1,p_2,\ldots,p_r\}$: conjunto de profesores.
  \item $S = \{s_1,s_2,\ldots,s_q\}$: conjunto de salones.
  \item $D = \{1,2,3,4,5\}$: conjunto ordenado de días académicos de la
        semana.
  \item $T$: conjunto ordenado de instantes de tiempo considerados durante
        un día académico.
  \item $F = \{f_1,f_2,\ldots,f_t\}$: conjunto de franjas horarias
        disponibles durante la semana.
\end{itemize}

Cada franja $f\in F$ se representa mediante una terna:

\[
f=\bigl(d(f),t_{\mathrm{inicio}}(f),t_{\mathrm{fin}}(f)\bigr)
\in D\times T\times T,
\]
con $t_{\mathrm{inicio}}(f)<t_{\mathrm{fin}}(f)$. Esta representación permite
ordenar temporalmente las clases, calcular los tiempos libres y determinar
cuándo dos franjas se superponen.

Se define la relación de solapamiento $\bowtie$ sobre $F$. Para
$f_i,f_j\in F$:

\[
f_i\bowtie f_j
\iff
d(f_i)=d(f_j)
\land
t_{\mathrm{inicio}}(f_i)<t_{\mathrm{fin}}(f_j)
\land
t_{\mathrm{inicio}}(f_j)<t_{\mathrm{fin}}(f_i).
\]

Dos clases son consecutivas cuando ocurren el mismo día, no existe otra clase
asignada entre ellas y la hora de inicio de la segunda no es anterior a la
hora de finalización de la primera. Esta relación se utilizará posteriormente
para calcular tiempos libres y distancias entre salones consecutivos.

\subsection{Funciones sobre los grupos}

Cada grupo $g\in G$ pertenece exactamente a una asignatura. Se definen las
siguientes funciones:

\[
\begin{aligned}
  \operatorname{asignatura} &: G \rightarrow A,\\
  \operatorname{profesor}   &: G \rightarrow P,\\
  \operatorname{horario}    &: G \rightarrow 2^F,\\
  \operatorname{salon}      &: G \rightarrow S,\\
  \operatorname{capacidad}  &: G \rightarrow \mathbb{N}.
\end{aligned}
\]

Para cada grupo $g$, $\operatorname{asignatura}(g)$ indica la asignatura a la
que pertenece; $\operatorname{profesor}(g)$, el profesor encargado;
$\operatorname{horario}(g)$, el conjunto de franjas que ocupa;
$\operatorname{salon}(g)$, el salón donde se desarrolla; y
$\operatorname{capacidad}(g)$, el número máximo de estudiantes que puede
recibir.

\subsection{Información de los estudiantes}

Para cada estudiante $e\in E$ se define el conjunto de asignaturas solicitadas:

\[
\operatorname{Solicitudes}(e)\subseteq A.
\]

También se define la función de prioridad:

\[
\operatorname{prioridad}:E\times A\rightarrow\{1,2,\ldots,K\},
\]

donde $\operatorname{prioridad}(e,a)$ representa la importancia que tiene para
el estudiante $e$ cursar la asignatura $a$. Un valor mayor indica una prioridad
mayor.

\subsection{Información de los grupos}

Para cada grupo $g\in G$ se conoce:

\[
\operatorname{asignatura}(g)\in A,\quad
\operatorname{profesor}(g)\in P,\quad
\operatorname{horario}(g)\subseteq F,
\]

\[
\operatorname{salon}(g)\in S,\qquad
\operatorname{capacidad}(g)\in\mathbb{N}.
\]

\section{Preferencias de los estudiantes}

\subsection{Preferencia por profesor}

Se define la función:

\[
\operatorname{preferenciaProfesor}:E\times P
\rightarrow\{0,1,2,3,4,5\},
\]

donde $\operatorname{preferenciaProfesor}(e,p)$ representa el nivel de
preferencia del estudiante $e$ por recibir clase con el profesor $p$, según la
siguiente escala:

\begin{table}[H]
\centering
\caption{Escala de preferencia por profesor.}
\begin{tabular}{cl}
\toprule
\textbf{Valor} & \textbf{Significado} \\
\midrule
0 & Sin preferencia o profesor desconocido \\
1 & Preferencia muy baja \\
2 & Preferencia baja \\
3 & Preferencia media \\
4 & Preferencia alta \\
5 & Preferencia muy alta \\
\bottomrule
\end{tabular}
\end{table}

\subsection{Preferencias de horario}

Para cada estudiante $e\in E$ se define:

\[
\operatorname{PreferenciasHorario}(e)=
\left(h_{\min}(e),h_{\max}(e),c(e),l(e)\right),
\]

donde:

\begin{itemize}
  \item $h_{\min}(e)\in T$ es la hora mínima deseada para iniciar clases;
  \item $h_{\max}(e)\in T$ es la hora máxima deseada para terminar clases;
  \item $c(e)\in\{0,1\}$ expresa la preferencia por un horario compacto
        ($1$ si prefiere evitar huecos y $0$ en caso contrario);
  \item $l(e)\subseteq D$ es el conjunto de días de la semana en
        los que el estudiante prefiere no tener clases.
\end{itemize}

\subsection{Cercanía entre salones}

Se define la función:

\[
\operatorname{distancia}:S\times S\rightarrow\mathbb{R}_{\geq 0},
\]

donde $\operatorname{distancia}(s_i,s_j)$ representa la distancia física en
metros entre los salones $s_i$ y $s_j$. Para una instancia real se pueden
utilizar mediciones del campus de la Pontificia Universidad Javeriana.
En los experimentos se utilizan distancias sintéticas.

\subsection{Preferencias sociales y académicas}

Para cada estudiante $e\in E$ se define:

\[
\operatorname{Compañeros}(e)\subseteq E\setminus\{e\},
\]

que representa el conjunto de estudiantes con los cuales $e$ preferiría
coincidir en uno o más grupos.

\section{Precondiciones}

Las entradas deben satisfacer las siguientes propiedades antes de ejecutar
cualquier algoritmo de asignación:

\begin{enumerate}
  \item Los conjuntos $E$, $A$, $G$, $P$, $S$, $D$, $T$ y $F$ son finitos.
  \item Para toda franja $f\in F$:
        \[
        d(f)\in D,\qquad
        t_{\mathrm{inicio}}(f),t_{\mathrm{fin}}(f)\in T,
        \]
        \[
        t_{\mathrm{inicio}}(f)<t_{\mathrm{fin}}(f).
        \]
  \item Para todo $g\in G$:
        \[
        \operatorname{asignatura}(g)\in A,\quad
        \operatorname{profesor}(g)\in P,\quad
        \operatorname{horario}(g)\subseteq F,
        \]
        \[
        \operatorname{salon}(g)\in S,\qquad
        \operatorname{capacidad}(g)>0.
        \]
  \item Para todo $e\in E$:
        \[
        \emptyset\neq\operatorname{Solicitudes}(e)\subseteq A.
        \]
  \item Para todo par $s_i,s_j\in S$:
        \[
        \operatorname{distancia}(s_i,s_j)\geq 0,
        \]
        \[
        \operatorname{distancia}(s_i,s_i)=0,
        \]
        \[
        \operatorname{distancia}(s_i,s_j)
        =\operatorname{distancia}(s_j,s_i).
        \]
  \item Para todo $e\in E$:
        \[
        \operatorname{Compañeros}(e)\subseteq E\setminus\{e\}.
        \]
\end{enumerate}

\section{Salida formal y postcondiciones}

La salida principal del problema es la función:

\[
H:E\rightarrow 2^G,
\]

donde $H(e)$ representa el conjunto de grupos asignados al estudiante $e$.
La salida completa puede expresarse como:

\[
H=\left\{\bigl(e,H(e)\bigr)\mid e\in E\right\}.
\]

Cada conjunto $H(e)$ debe cumplir las siguientes postcondiciones.

\begin{restriccion}[Asignación de asignaturas solicitadas]
Para todo estudiante $e\in E$ y todo grupo $g\in H(e)$:
\[
\operatorname{asignatura}(g)\in\operatorname{Solicitudes}(e).
\]
\end{restriccion}

\begin{restriccion}[Ausencia de cruces de horario]
Para todo estudiante $e\in E$, todo par de grupos distintos
$g_i,g_j\in H(e)$ y toda pareja de franjas
$f_i\in\operatorname{horario}(g_i)$,
$f_j\in\operatorname{horario}(g_j)$:
\[
\neg\bigl(f_i\bowtie f_j\bigr).
\]
\end{restriccion}

\begin{restriccion}[Una sola sección por asignatura]
Para todo estudiante $e\in E$ y todo par de grupos distintos
$g_i,g_j\in H(e)$:
\[
\operatorname{asignatura}(g_i)\neq\operatorname{asignatura}(g_j).
\]
\end{restriccion}

\begin{restriccion}[Respeto de la capacidad]
Para todo grupo $g\in G$:
\[
\left|\{e\in E\mid g\in H(e)\}\right|
\leq\operatorname{capacidad}(g).
\]
\end{restriccion}

\section{Métricas de calidad}

Una vez comprobado que $H$ satisface las postcondiciones, se asocia a cada
estudiante un vector de métricas:

\[
M_e=\left(m_1(e),m_2(e),\ldots,m_z(e)\right).
\]

La salida extendida se define como:

\[
\operatorname{Resultado}(e)=\bigl(H(e),M_e\bigr).
\]

Las métricas pueden incluir:

\begin{itemize}
  \item número de asignaturas solicitadas y asignadas;
  \item satisfacción de las asignaturas de mayor prioridad;
  \item cantidad total de horas libres entre clases;
  \item distancia total recorrida entre clases consecutivas;
  \item satisfacción promedio con los profesores asignados;
  \item coincidencias con compañeros preferidos.
\end{itemize}

Para formalizar la cobertura, se define primero el conjunto de asignaturas
asignadas al estudiante $e$:

\[
\operatorname{Asignadas}(e)=
\left\{a\in\operatorname{Solicitudes}(e)\,\middle|\,
\exists g\in H(e):\operatorname{asignatura}(g)=a\right\}.
\]

La cobertura individual es:

\[
\operatorname{Cobertura}(e)=
\frac{|\operatorname{Asignadas}(e)|}
     {|\operatorname{Solicitudes}(e)|}.
\]

La cobertura global es:

\[
\operatorname{CoberturaGlobal}=
\frac{\displaystyle\sum_{e\in E}|\operatorname{Asignadas}(e)|}
     {\displaystyle\sum_{e\in E}|\operatorname{Solicitudes}(e)|}.
\]

Estas métricas permiten comparar asignaciones válidas. El orden en que se
utilizan para elegir entre dos resultados se establece en la
Sección~\ref{sec:criterio-final}.

\subsection{Ejemplo conceptual de salida}

\begin{table}[H]
\centering
\caption{Ejemplo de horario asignado al estudiante $e_1$.}
\begin{tabular}{ll}
\toprule
\textbf{Asignatura} & \textbf{Grupo asignado} \\
\midrule
Análisis de Algoritmos & G03 \\
Bases de Datos         & G07 \\
Computación Móvil      & G11 \\
Optimización           & G15 \\
\bottomrule
\end{tabular}
\end{table}

\begin{table}[H]
\centering
\caption{Métricas del horario asignado al estudiante $e_1$.}
\begin{tabular}{lr}
\toprule
\textbf{Métrica} & \textbf{Valor} \\
\midrule
Materias solicitadas        & 4 \\
Materias asignadas          & 4 \\
Cruces de horario           & 0 \\
Tiempo libre total          & 6 horas \\
Distancia total recorrida   & 850 metros \\
Satisfacción con profesores & 4.25 / 5 \\
Prioridades satisfechas     & 100\% \\
\bottomrule
\end{tabular}
\end{table}

\part{Estrategia de datos}

\section{Generación de datos sintéticos}

El generador de datos sintéticos se controla mediante los siguientes
parámetros:

\begin{itemize}
  \item número de estudiantes;
  \item número de asignaturas;
  \item número de grupos por asignatura;
  \item número de profesores;
  \item número de salones;
  \item número de franjas horarias;
  \item capacidad promedio de los grupos.
\end{itemize}

A partir de estos parámetros se generan estudiantes, asignaturas, grupos,
profesores, horarios, salones, capacidades, prioridades, preferencias y
relaciones entre estudiantes. La generación parte de una semilla. En los
experimentos se utilizaron las semillas 202630, 202631 y 202632.

El generador crea primero los catálogos y los grupos. Cada grupo tiene
una o dos sesiones en días distintos, con duraciones de 60, 90 o 120
minutos. Los inicios se eligen entre las 08:00 y las 16:00 en pasos de
30 minutos. Si $R$ es el total de solicitudes, el cupo base se calcula
como $b=\max(1,\lfloor 0{,}8R/|G|\rfloor)$. Para cada grupo se aplica
una variación de hasta $\lfloor b/3\rfloor$, con cupo mínimo de uno.

Cada estudiante recibe un conjunto de asignaturas solicitadas sin
duplicados. Se usan pesos distintos para simular una demanda desigual
y se asigna una prioridad de 1 a 5 a cada solicitud. Las distancias
sintéticas se obtienen a partir de coordenadas enteras de los salones
y se registran en metros.

\section{Generación de salidas esperadas}

La estrategia de pruebas debe incluir no solamente los datos de entrada, sino
también resultados de referencia que permitan determinar si la salida de una
implementación es correcta. Debido a que un mismo caso puede admitir múltiples
horarios válidos, una salida esperada no se limita necesariamente a una única
asignación exacta.

Para cada caso de prueba se almacenarán:

\begin{itemize}
  \item los datos de entrada del caso;
  \item una o más asignaciones válidas de referencia, cuando sea posible
        determinarlas;
  \item los valores esperados de las métricas de cobertura y calidad;
  \item las propiedades de validez que toda salida aceptable debe satisfacer;
  \item la indicación de si el caso admite o no una asignación completa.
\end{itemize}

En los casos pequeños, las salidas de referencia se comprueban mediante
revisión manual y enumeración exhaustiva. La enumeración permite verificar
si existe una asignación completa y calcular el mejor resultado según el
criterio de comparación.

También se podrán construir casos con una \emph{solución sembrada}. Para ello,
primero se define una asignación de referencia $H^{*}$ y después se generan las
solicitudes, capacidades y franjas horarias de forma que $H^{*}$ satisfaga las
postcondiciones. A continuación pueden introducirse conflictos controlados,
como grupos con cupos reducidos o franjas solapadas. Cada caso se verifica
de forma independiente para conservar las propiedades esperadas.

Las salidas producidas por cualquier implementación serán evaluadas mediante
un verificador independiente. Este comprobará las cuatro postcondiciones,
recalculará las métricas y comparará sus valores con los resultados de
referencia. En escenarios medianos y grandes, donde enumerar todas las
soluciones no es viable, la salida esperada se expresará mediante condiciones
verificables, cotas conocidas y soluciones base, en lugar de exigir que el
algoritmo reproduzca exactamente un horario particular.

El archivo \texttt{salidas\_esperadas.json} registra, para cada caso,
la terna esperada para el voraz, el backtracking y la fuerza bruta. También
indica si se pueden atender todas las solicitudes.

\section{Escenarios de prueba}

\subsection{Escenario pequeño}

\begin{itemize}
  \item 5 estudiantes.
  \item 5 asignaturas.
  \item 2 grupos por asignatura.
\end{itemize}

Este escenario permitirá verificar manualmente la validez de los resultados.

\subsection{Escenario mediano}

\begin{itemize}
  \item 50 estudiantes.
  \item 15 asignaturas.
  \item Entre 3 y 4 grupos por asignatura.
\end{itemize}

Este escenario permitirá evaluar el comportamiento del algoritmo en
condiciones más cercanas a un entorno real.

\subsection{Escenario grande}

\begin{itemize}
  \item 500 o más estudiantes.
  \item 30 o más asignaturas.
  \item Múltiples grupos por asignatura.
\end{itemize}

Este escenario permitirá evaluar la escalabilidad y el tiempo de ejecución del
algoritmo.

\section{Escenarios especiales}

Se propone generar casos deliberadamente conflictivos, entre ellos:

\begin{itemize}
  \item muchas asignaturas con horarios coincidentes;
  \item grupos con baja capacidad;
  \item gran cantidad de estudiantes interesados en el mismo profesor;
  \item asignaturas con una única sección disponible;
  \item estudiantes con muchas asignaturas de alta prioridad;
  \item salones físicamente muy separados;
  \item solicitudes imposibles de satisfacer completamente.
\end{itemize}

Estos casos permitirán evaluar la robustez del algoritmo.

\section{Datos cercanos a escenarios reales}

Como segunda etapa se propone utilizar información obtenida de fuentes reales
o construida a partir de encuestas. Los datos podrían provenir de horarios
públicos de asignaturas, información suministrada voluntariamente por
estudiantes, datos anonimizados e información académica disponible
públicamente.

Las preferencias subjetivas se obtendrán mediante un formulario con preguntas
como:

\begin{itemize}
  \item ¿Qué asignaturas necesita cursar?
  \item ¿Qué prioridad tiene cada asignatura?
  \item ¿Tiene profesores preferidos?
  \item ¿Cuál es la hora mínima en la que desea iniciar clases?
  \item ¿Cuál es la hora máxima en la que desea terminar clases?
  \item ¿Prefiere horarios compactos?
  \item ¿Desea tener un día libre?
  \item ¿Con qué estudiantes desea coincidir?
\end{itemize}

No se recolectará información personal sensible que no sea necesaria para el
problema.

\section{Formato de almacenamiento}
\label{sec:csv}

Cada instancia se almacena en una carpeta con once archivos CSV. Una fila
representa una entidad o una relación. Por ejemplo, si un grupo se reúne
dos días, sus datos fijos se escriben una sola vez en \texttt{grupos.csv}
y las dos sesiones se escriben como filas en \texttt{horarios.csv}.
Así se pueden añadir sesiones o solicitudes sin modificar el encabezado
ni repetir el cupo, el profesor o el salón de un grupo.

\subsection{Convenciones y diccionario de datos}

Todos los CSV tienen encabezado obligatorio, columnas en el orden documentado,
separador coma y codificación UTF-8. Las columnas de un registro deben estar
completas. Los identificadores se tratan como texto y deben conservarse
al importar o editar los archivos. Los campos numéricos son enteros decimales
sin relleno con ceros. Las horas se expresan en minutos desde medianoche
(\texttt{480} equivale a 08:00). Los días van de 1 (lunes) a 5 (viernes).
Las métricas y los grupos asignados se guardan como salida JSON en una ruta
separada de las entradas.

\begin{table}[H]
\centering
\caption{Catálogos y entidades principales del esquema CSV.}
\small
\begin{tabular}{@{}p{0.27\linewidth}p{0.69\linewidth}@{}}
\toprule
\textbf{Archivo} & \textbf{Columnas en orden y clave primaria} \\
\midrule
\texttt{asignaturas.csv} &
\texttt{asignatura\_id}. Clave: asignatura. \\
\texttt{profesores.csv} &
\texttt{profesor\_id}. Clave: profesor. \\
\texttt{salones.csv} &
\texttt{salon\_id}. Clave: salón. \\
\texttt{estudiantes.csv} &
\texttt{estudiante\_id, inicio\_min, fin\_max, compacto}.\newline
Clave: estudiante. \texttt{compacto} vale 0 o 1. \\
\texttt{grupos.csv} &
\texttt{grupo\_id, asignatura\_id, profesor\_id,}\newline
\texttt{salon\_id, capacidad}. Clave: grupo. Asignatura,
profesor y salón deben existir en sus catálogos. \\
\bottomrule
\end{tabular}
\end{table}

\begin{table}[H]
\centering
\caption{Sesiones, solicitudes, preferencias y distancias.}
\small
\begin{tabular}{@{}p{0.27\linewidth}p{0.69\linewidth}@{}}
\toprule
\textbf{Archivo} & \textbf{Columnas en orden y clave primaria} \\
\midrule
\texttt{horarios.csv} &
\texttt{grupo\_id, dia, inicio, fin}. Clave: los cuatro campos.
El grupo debe existir en \texttt{grupos.csv}. \\
\texttt{solicitudes.csv} &
\texttt{estudiante\_id, asignatura\_id, prioridad}. Clave:
(estudiante, asignatura). Ambos identificadores deben existir. \\
\texttt{preferencias\_}\newline\texttt{profesor.csv} &
\texttt{estudiante\_id, profesor\_id, valor}. Clave:
(estudiante, profesor). Valor entre 0 y 5. \\
\texttt{dias\_libres.csv} &
\texttt{estudiante\_id, dia}. Clave: (estudiante, día). \\
\texttt{companeros.csv} &
\texttt{estudiante\_id, companero\_id}. Clave: (estudiante,
compañero). Ambos son estudiantes distintos. La relación es dirigida. \\
\texttt{distancias.csv} &
\texttt{origen\_id, destino\_id, metros}. Clave: (origen, destino).
Se registra cada par ordenado de salones, incluida la diagonal. \\
\bottomrule
\end{tabular}
\end{table}

Las tablas de preferencias por profesor, días libres y compañeros pueden
contener solo el encabezado. La ausencia de una fila indica que no se
informó esa preferencia. Si un estudiante no define límites de hora, la
exportación registra \texttt{inicio\_min=0}, \texttt{fin\_max=1440}.
Cada estudiante tiene al menos una solicitud y cada grupo al menos una
sesión. Una asignatura puede estar en el catálogo y en una solicitud aunque
carezca de grupos. Esa solicitud podrá quedar pendiente.

\subsection{Ejemplo de filas enlazadas}

Estos fragmentos muestran cómo se reconstruye una parte de la instancia.
El estudiante \texttt{E001} solicita \texttt{A01} con prioridad 5.
El grupo \texttt{A01-G1} pertenece a esa asignatura, admite 12 estudiantes y
ocupa dos sesiones en el mismo salón. El lector reúne las dos franjas bajo
el mismo identificador.

\begin{verbatim}
estudiantes.csv
estudiante_id,inicio_min,fin_max,compacto
E001,480,1080,1

solicitudes.csv
estudiante_id,asignatura_id,prioridad
E001,A01,5

grupos.csv
grupo_id,asignatura_id,profesor_id,salon_id,capacidad
A01-G1,A01,P01,S01,12

horarios.csv
grupo_id,dia,inicio,fin
A01-G1,1,480,570
A01-G1,3,480,570
\end{verbatim}

Los catálogos deben contener \texttt{A01}, \texttt{P01} y \texttt{S01}.
Si la instancia solo tiene el salón \texttt{S01}, \texttt{distancias.csv}
contiene el encabezado \texttt{origen\_id,destino\_id,metros} y la fila
\texttt{S01,S01,0}.

\subsection{Creación, carga y mantenimiento}

\begin{enumerate}
  \item Crear los catálogos con identificadores estables. Para una fuente
        real, convertir horarios a minutos y reemplazar identificadores
        personales por códigos de la instancia antes de guardar.
  \item Crear los grupos con referencias a los catálogos y registrar sus
        franjas en filas separadas. Crear estudiantes, solicitudes y
        preferencias usando esos mismos identificadores.
  \item Ordenar las filas por clave y escribir los once CSV en una carpeta
        nueva. Para datos sintéticos, registrar semilla, escenario, versión
        del esquema y versión de Python en \texttt{meta.json}.
  \item Leer primero los catálogos y estudiantes. Comprobar después las
        referencias al cargar grupos, sesiones, solicitudes, preferencias
        y distancias. Reconstruir una lista de sesiones por grupo y un
        conjunto de solicitudes y preferencias por estudiante.
  \item Validar la instancia reconstruida antes de asignar. Al corregir o
        regenerar una instancia, guardar el conjunto como otra carpeta.
        Así se conservan los datos usados en las mediciones anteriores.
\end{enumerate}

La interfaz permite generar la carpeta mediante \textbf{Generar CSV},
abrirla mediante \textbf{Abrir carpeta CSV} y revisar cada archivo en la
pestaña \textbf{Datos CSV}. Para automatizar las pruebas, se puede reproducir
el mismo flujo desde la carpeta \texttt{implementacion}.

\begin{verbatim}
python horarios.py generar pequeno ejemplo --formato csv --semilla 202630
python horarios.py resolver ejemplo resultado.json
python horarios.py verificar ejemplo resultado.json
\end{verbatim}

\section{Estrategia de validación}

Antes de asignar se comprueba la estructura de los CSV, la unicidad
de sus claves y la existencia de los identificadores relacionados.
También se revisan los rangos de las horas, la capacidad de los grupos
y la coherencia de las distancias.

La validación de las salidas se divide en dos niveles.

\subsection{Nivel 1: validez}

Se verifica que no existan cruces de horario, que no se exceda la capacidad de
ningún grupo, que cada estudiante tenga máximo una sección por asignatura y que
los grupos asignados correspondan a asignaturas solicitadas. Si alguna
condición falla, la solución se considera inválida.

\subsection{Nivel 2: calidad}

Entre soluciones válidas se comparan métricas de calidad, entre ellas la
cobertura individual, la satisfacción de prioridades, el tiempo libre total,
la distancia recorrida, la satisfacción con profesores y la coincidencia con
compañeros preferidos.

Para el conjunto completo de estudiantes se evalúan las siguientes métricas
globales:

\begin{itemize}
  \item cobertura global de asignaciones;
  \item satisfacción promedio de prioridades;
  \item tiempo libre promedio entre clases;
  \item distancia promedio recorrida;
  \item satisfacción promedio con los profesores;
  \item tiempo de ejecución del algoritmo.
\end{itemize}

\section{Alcance de esta entrega}

La definición del problema y la estrategia de datos establecen el
contrato de entrada y salida. A partir de este contrato, la parte
siguiente presenta el algoritmo implementado, el análisis de su costo
y la evaluación de los horarios obtenidos.

\clearpage
\part{Implementación y resultados}

\section{Algoritmo de asignación}

Al leer los CSV, el programa reúne las solicitudes, los grupos y las
franjas en la instancia \texttt{x}. A partir de esos datos se construye
la asignación \(H\), se comprueba si cumple las restricciones y se
calculan las métricas. Tanto la interfaz como la terminal usan esas
mismas funciones de Python. Los tres algoritmos están en
\texttt{implementacion/algoritmo.py}.

\subsection{Criterio de comparación}
\label{sec:criterio-final}

Para cada asignación válida se compara la terna \(F(H)=\bigl(N(H),W(H),Q(H)\bigr)\)
en orden lexicográfico. Los dos primeros componentes se definen mediante
las solicitudes atendidas y sus prioridades.
\[
\begin{aligned}
N(H)&=\sum_{e\in E}\left|\operatorname{Asignadas}(e)\right|,\\
W(H)&=\sum_{e\in E}\ \sum_{a\in\operatorname{Asignadas}(e)}
        \operatorname{prioridad}(e,a).
\end{aligned}
\]
Si dos asignaciones empatan en \(N(H)\), se compara \(W(H)\).
El puntaje de preferencias \(Q(H)\) se utiliza cuando también
coinciden las prioridades atendidas.

Para calcular \(Q(H)\) se suman los puntajes individuales. Para cada
estudiante se suman 60 puntos por unidad de preferencia del profesor
asignado y 120 por cada coincidencia de un compañero preferido en un
grupo. Se descuenta un punto por minuto de clase fuera de la ventana
deseada y 120 por sesión programada en un día libre preferido. Para
los estudiantes que prefieren un horario compacto se resta un punto
por minuto de hueco. También se resta la parte entera de la distancia
total recorrida en metros dividida entre diez. Los huecos y las
distancias se miden entre sesiones consecutivas del mismo día.

\subsection{Línea base voraz}

La función \texttt{construir\_voraz(x, alternativas, cruzan)} visita
las solicitudes en orden de identificador de estudiante y asignatura.
Para cada solicitud toma el primer grupo compatible que tenga cupo.
Si no lo encuentra, deja la solicitud pendiente. Esta asignación es
el punto de partida del backtracking.

\subsection{Backtracking con poda}

La función \texttt{backtracking(x, limite\_nodos)} empieza por las
solicitudes que tienen menos grupos. Si dos tienen la misma cantidad
de grupos, va primero la de mayor prioridad. Si persiste el empate,
se usan los identificadores.
En cada nivel se prueba un grupo que tenga cupo y que no se cruce con
el horario parcial del estudiante. También se puede dejar la materia
pendiente. Para probar otra opción, se deshace la última decisión y se
devuelve el cupo si se había usado. La pila guarda el punto del árbol
desde el que se retoma la búsqueda.

Para acotar lo que puede aportar una rama, sean \(D_a\) el número de
solicitudes pendientes de explorar para la asignatura \(a\), y \(C_a\)
el total de cupos libres de sus grupos. En ese estado se pueden atender
como máximo \(u_a=\min(D_a,C_a)\) solicitudes adicionales de \(a\).
La suma de las \(u_a\) da una cota de cobertura. Para la prioridad se
suman los \(u_a\) valores más altos de las solicitudes restantes de cada
materia. Como no se tienen en cuenta los cruces futuros al calcular la
cota, esta puede sobreestimar el resultado alcanzable. Se descarta la
rama cuando ni siquiera esa cobertura y prioridad pueden igualar las
dos primeras componentes del mejor resultado conocido. Si la cota
iguala ambas, se sigue explorando porque el puntaje de preferencias
aún puede mejorar.

\subsection{Pseudocódigo}

El siguiente pseudocódigo usa llamadas recursivas para describir el
recorrido. En el programa, cada nivel equivale a una entrada de la pila.

\begin{verbatim}
mejor <- resultado de la linea base voraz
Explorar(i, estado):
    si se alcanzo el limite:
        registrar busqueda incompleta y detener el recorrido
    contar este nodo
    si la cota de cobertura y prioridad es inferior a mejor:
        retornar
    si i es el numero de solicitudes:
        comparar estado con mejor y actualizar si mejora
        retornar
    para cada grupo compatible con cupo y la opcion pendiente:
        aplicar la decision de la solicitud i
        Explorar(i + 1, estado)
        deshacer la decision y restituir los cupos

devolver mejor y registrar si se agoto la busqueda
\end{verbatim}

El límite predeterminado es de 15\,000 nodos. Al alcanzarlo con ramas
todavía por explorar, el programa conserva la mejor asignación válida
encontrada e indica que se alcanzó el límite. Si se agotan las ramas
sin superar ese presupuesto, el programa informa que se comprobó el
óptimo bajo \(F(H)\).

\subsection{Correctitud y terminación}

En cada nivel, todo grupo asignado pertenece a una solicitud ya
procesada, se asigna a lo sumo una sección por materia y las sesiones
del estudiante son compatibles. El cupo libre de cada grupo es
su capacidad menos el número de estudiantes asignados. El estado
vacío cumple estas condiciones. Al avanzar, solo se asignan grupos
compatibles con cupo y se actualiza el contador. Al retroceder se
restaura el estado anterior. La opción pendiente conserva la validez
del horario y consume la solicitud en la rama explorada.

Cada asignación válida corresponde a una secuencia de decisiones
de grupo u omisión. La poda descarta únicamente ramas cuya cota
queda por debajo del mejor par de cobertura y prioridad. Cuando se
agota el árbol, la mejor asignación es óptima según \(F(H)\). Si se
alcanza el límite, la salida sigue siendo válida y tiene un valor
al menos igual al de la base, que fue el resultado inicial.

\subsection{Fuerza bruta}

La función \texttt{fuerza\_bruta(x, limite)} busca el óptimo por
enumeración exhaustiva. Recorre todas las combinaciones de grupos y
omisiones. Para cada combinación arma un horario y comprueba si
respeta las restricciones. Entre los horarios válidos conserva el
de mayor \(F(H)\). En los casos pequeños se usa como referencia para
comprobar el resultado del backtracking.

\section{Complejidad de los algoritmos}
\label{sec:complejidad-algoritmos}

Sea \(R\) el total de solicitudes y \(r\) el número de solicitudes que
tienen al menos un grupo. Llamamos \(B\) al máximo de grupos por
asignatura, \(K\) al máximo de solicitudes por estudiante y \(L\) al
máximo de sesiones por grupo. Comparar dos grupos por primera vez
cuesta hasta \(O(L^2)\), pues se compara cada sesión de uno con las del
otro. El análisis parte de los datos de los CSV ya cargados en memoria
y cuenta el espacio adicional que usa cada método.

\subsection{Línea base voraz}

Para cada solicitud, el voraz revisa hasta \(B\) grupos. Por cada grupo
con cupo, comprueba que sea compatible con los grupos que el estudiante
ya tiene, que son a lo sumo \(K\). En el peor caso, esa revisión cuesta
\(O(RBKL^2)\). Al sumar el costo de ordenar estudiantes, grupos y
solicitudes, se obtiene
\begin{multline*}
 T_{\mathrm{base}} = O\bigl(|E|\log(2+|E|)+|G|\log(2+|G|)\\
                          {}+ R\log(2+K)+RBKL^2\bigr).
\end{multline*}
Las listas de grupos, los cupos y los horarios ocupan
\(O(|E|+|G|+R)\). El programa también guarda hasta
\(\min(|G|^2,100\,000)\) resultados de cruces para reutilizarlos.

\subsection{Backtracking con poda}

Cada nivel del árbol corresponde a una solicitud. Se puede asignar
uno de sus grupos o dejarla pendiente. Con \(r\) niveles y hasta
\(B+1\) opciones por nivel, el árbol puede tener \(O((B+1)^r)\) nodos.
La poda corta una rama cuando, aun en el mejor caso, sus decisiones
pendientes no bastarían para mejorar el horario guardado.

Sea \(V\) el número de nodos visitados y \(U\) el de hojas que se
evalúan. El método parte del horario voraz, revisa hasta \(B\) grupos
en cada nodo y llama a \texttt{evaluar()} en esas \(U\) hojas. Si
\(T_F\) es el costo de esa función, el tiempo se acota por
\begin{multline*}
 T_{\mathrm{bt}} =
 O\bigl(T_{\mathrm{base}}+T_F+R\log(2+R)\\
         {}+V(1+BKL^2)+UT_F\bigr).
\end{multline*}
El término \(R\log(2+R)\) corresponde a ordenar los pedidos antes de
recorrer el árbol. El límite de 15\,000 nodos acota \(V\) en la
configuración usada en las pruebas.

La pila tiene hasta \(r\) niveles. El espacio adicional es
\(O(|E|+|G|+R+KL+\min(|G|^2,100\,000))\). Ahí se cuentan el horario
parcial, el mejor horario, los cupos, las sesiones que se ordenan al
evaluar y la caché de cruces.

\subsection{Fuerza bruta}

Si el pedido \(i\) tiene \(b_i\) grupos, hay \(1+b_i\) opciones
contando la posibilidad de dejarlo pendiente. En el código,
\texttt{combinaciones} guarda el producto
\[
 J=\prod_{i=1}^{r}(1+b_i)\le (B+1)^r.
\]
Con dos pedidos de dos grupos cada uno, por ejemplo, salen nueve
combinaciones. Para cada una se construye un horario candidato y se
comprueba su validez. Los candidatos válidos se evalúan. Si \(T_V\)
es el costo de \texttt{validar()} y \(T_{\mathrm{prep}}\) el de
preparar las opciones, el tiempo en el peor caso es
\[
 O\bigl(T_{\mathrm{prep}}+J(|E|\log(2+|E|)+r+T_V+T_F)\bigr).
\]
En el código se ordenan los estudiantes al crear cada candidato.
La preparación incluye ordenar los grupos y los pedidos y reunir
las opciones de cada solicitud:
\begin{multline*}
 T_{\mathrm{prep}}=O\biggl(|E|\log(2+|E|)+|G|\log(2+|G|)\\
                    {}+R\log(2+K)+\sum_{i=1}^{r}(1+b_i)\biggr).
\end{multline*}
La enumeración sirve como referencia para instancias de hasta
50\,000 combinaciones. Genera una combinación a la vez, así que el
espacio adicional es \(O(|E|+|G|+R+KL+\sum_{i=1}^{r}(1+b_i))\).
Se guardan las listas de opciones, el candidato que se revisa y el
mejor horario encontrado.

Al validar se ordenan hasta \(KL\) sesiones por estudiante y se
recuentan los cupos. Al evaluar, además, se revisan las coincidencias
con compañeros y se ordenan las materias pendientes. Si \(C\) es el
máximo de compañeros preferidos por estudiante, se tiene
\begin{align*}
 T_V &= O\bigl(|E|+R+|E|KL\log(2+KL)+|G|\log(2+|G|)\bigr),\\
 T_F &= O\bigl(|E|\log(2+|E|)
       + |E|\bigl(KL\log(2+KL)+CK^2+K\log(2+K)\bigr)\bigr).
\end{align*}

\section{Verificación y evaluación experimental}

El verificador reconstruye los cupos ocupados y ordena las sesiones por
estudiante para detectar cruces. También revisa que la salida incluya
exactamente a los estudiantes de la entrada, que las secciones existan
y que correspondan a solicitudes sin repetir materias. Este procedimiento
lee la asignación final y calcula la validez sin depender del estado interno
de cupos del solucionador.

\subsection{Casos de referencia}

Los casos de referencia permiten comprobar el resultado en situaciones que
se pueden resolver a mano. La tabla compara la cantidad de solicitudes
atendidas para cada método.

\begin{table}[H]
\centering
\caption{Solicitudes atendidas en los casos de referencia.}
\begin{tabular}{lccc}
\toprule
Caso & Base voraz & Backtracking & Fuerza bruta \\
\midrule
Cambio de sección & 1 de 2 & 2 de 2 & 2 de 2 \\
Cupo compartido & 1 de 2 & 1 de 2 & 1 de 2 \\
Reasignación entre estudiantes & 2 de 4 & 3 de 4 & 3 de 4 \\
\bottomrule
\end{tabular}
\end{table}

En el caso de cambio de sección, A1 ocupa el lunes de 08:00 a 09:00 y
bloquea las dos opciones de B. El backtracking retrocede y prueba A2,
de 09:00 a 10:00. Así puede asignar también B1. En el caso de cupo
compartido, la base atiende una solicitud de prioridad 1 y el
backtracking atiende la de prioridad 5 con la misma cobertura. En el
caso de reasignación, A1 queda libre para E2 al asignar A2 a E1.
El backtracking y la enumeración coinciden en las tres solicitudes
atendidas. La cuarta materia carece de grupos.

\subsection{Resultados por tamaño}

Cada valor de la tabla es el promedio de tres instancias generadas con las
semillas 202630, 202631 y 202632. Para cada instancia se usaron los mismos
datos con ambos métodos. El tiempo corresponde a la ejecución del
solucionador, sin incluir la lectura de archivos ni la validación externa.
Las mediciones se registraron en Windows con Python 3.14.8 y un procesador
Intel Core Ultra 7 255U.

\begin{table}[H]
\centering
\caption{Promedios de cobertura, prioridad atendida y tiempo de resolución.}
\small
\begin{tabular}{r rrr rrr}
\toprule
 & \multicolumn{3}{c}{Base voraz} & \multicolumn{3}{c}{Propuesto} \\
\cmidrule(lr){2-4}\cmidrule(lr){5-7}
Estudiantes & Cobertura & Prioridad & Tiempo & Cobertura & Prioridad & Tiempo \\
\midrule
5   & 62,22\% & 28,00   & 0,05 ms & 62,22\% & 33,00   & 4,01 ms \\
50  & 57,87\% & 429,67  & 0,26 ms & 58,13\% & 542,67  & 499,44 ms \\
500 & 65,87\% & 5953,00 & 3,62 ms & 65,87\% & 5953,00 & 48,30 ms \\
\bottomrule
\end{tabular}
\end{table}

En dos instancias medianas el backtracking atendió una solicitud más
que la base. En otras cuatro instancias mantuvo la cobertura y aumentó
la prioridad atendida. En las tres grandes se conservó la misma terna.
La búsqueda se agotó en las tres instancias pequeñas. En las seis
medianas y grandes se alcanzó el límite de 15\,000 nodos, por lo que
se reporta la mejor asignación encontrada en cada caso. Los tiempos
medidos dependen también de cuántas asignaciones completas se evaluaron.

En \repo{} están los resultados de cada semilla y el estado de cada
búsqueda. También se puede consultar cómo se probó la lectura de los
CSV, la validación de los horarios y el límite de nodos. En doce
instancias pequeñas se comparó el backtracking con la enumeración
exhaustiva.

\section{Ejecución para la sustentación}

La interfaz utiliza Python 3.10 o posterior con Tkinter. En PyCharm, se
abre la carpeta completa del repositorio \repo, se selecciona un
intérprete local y se ejecuta \texttt{implementacion/interfaz.py}
mediante \textbf{Run 'interfaz'}. En Windows también se puede abrir
\texttt{Abrir interfaz.pyw} con doble clic. La ventana carga el ejemplo
CSV al iniciar.

\begin{enumerate}
  \item Seleccionar el escenario pequeño, dejar la semilla 202630 y pulsar
        \textbf{Generar CSV}. En \textbf{Datos CSV}, seguir la solicitud de
        A01 de E001 desde \texttt{solicitudes.csv} hasta un grupo en
        \texttt{grupos.csv} y sus sesiones en \texttt{horarios.csv}.
  \item Pulsar \textbf{Asignar horarios}. Para esa entrada, el método
        propuesto atiende 10 de 15 solicitudes. Seleccionar E005 para
        mostrar su horario y la materia A04 pendiente.
  \item Pulsar \textbf{Verificar resultado} para mostrar el cumplimiento
        de las cuatro restricciones. El botón \textbf{Guardar resultado
        JSON} permite conservar la asignación y las métricas.
  \item En \textbf{Casos para la sustentación}, abrir \textbf{Cambio
        de sección}. Comparar \textbf{Base \(\cdot\) voraz} con
        \textbf{Propuesto \(\cdot\) backtracking}. El caso
        \textbf{Reasignación entre estudiantes} permite comparar
        el backtracking con \textbf{Fuerza bruta}. Al pie de la ventana
        aparece el número de nodos y el estado de la búsqueda.
\end{enumerate}

\section{Conclusión}

En seis de las nueve instancias, el backtracking atendió más
solicitudes que el voraz o logró una mayor suma de prioridades con la
misma cobertura. En las tres instancias grandes llegó al límite de
nodos y conservó el resultado del voraz. En las pequeñas pudo agotar
la búsqueda y comprobar el óptimo.

Con la interfaz se puede seguir cada solicitud desde el CSV hasta el
horario que recibe el estudiante. El resultado muestra qué materias
quedaron pendientes y permite comprobar si la asignación respeta los
cupos y evita los cruces. El estado de la búsqueda indica si se
comprobó el óptimo o si se llegó al límite de nodos.

\end{document}
